\section{Erd\H{o}s Problem 250: irrationality of a divisor-sum series}

\subsection{1) Statement and scope}
The number
\[
 S=\sum_{n=1}^{\infty}\frac{\sigma(n)}{2^n},
 \qquad \sigma(n)=\sum_{d\mid n}d,
\]
is not merely irrational: it is transcendental.
We give the complete deduction from Nesterenko's theorem on
Eisenstein series. That deep algebraic-independence theorem
is an external input and is not reproved here.

\subsection{2) Convergence and the modular-series identity}
For $r\geq1$, let $\sigma_r(n)=\sum_{d\mid n}d^r$.
The bound $\sigma_r(n)\leq n^{r+1}$ ensures convergence
of the following series for $|q|<1$:
\[
 \begin{split}
 P(q)&=1-24\sum_{n\geq1}\sigma_1(n)q^n,\\
 Q(q)&=1+240\sum_{n\geq1}\sigma_3(n)q^n,\\
 R(q)&=1-504\sum_{n\geq1}\sigma_5(n)q^n.
 \end{split}
\]
In particular,
\[
 P(1/2)=1-24S.
\]
The nonnegative divisor expansion also gives the Lambert-series
description
\[
 S=\sum_{d\geq1}\sum_{j\geq1}\frac{d}{2^{dj}}
  =\sum_{d\geq1}\frac{d}{2^d-1}.
\]
Both the interchange and the identity follow from absolute
convergence (or directly from Tonelli's theorem).

\subsection{3) Applying the external theorem}
Nesterenko's theorem states that for each complex $q$ with
$0<|q|<1$,
\[
 \operatorname{trdeg}_{\mathbb Q}
       \mathbb Q\bigl(q,P(q),Q(q),R(q)\bigr)\geq3.
\]
At $q=1/2$, adjoining $q$ adds no transcendence degree.
Thus the three numbers $P(1/2),Q(1/2),R(1/2)$ are algebraically
independent over $\mathbb Q$. In particular $P(1/2)$ is
transcendental: otherwise the field generated by all three
would have transcendence degree at most two.
If $S$ were algebraic, then $1-24S=P(1/2)$ would be algebraic,
a contradiction. This proves the asserted transcendence.

\subsection{4) Scope and sources}
Numerical approximation does not establish this conclusion.
For example the entirely elementary tail bound
\[
 0<S-\sum_{n=1}^M\frac{\sigma(n)}{2^n}
 \leq\sum_{n>M}\frac{n^2}{2^n}
 =\frac{M^2+4M+6}{2^M}
\]
certifies convergence but not irrationality. The transcendence
step above depends essentially on the stated theorem.

Sources: \url{https://www.erdosproblems.com/250};
Yu.~V.~Nesterenko, \emph{Modular functions and transcendence
questions}, Sbornik Math. 187 (1996), 1319--1348,
\url{https://doi.org/10.1070/SM1996v187n09ABEH000158};
W.~Zudilin, \emph{Diophantine problems for $q$-zeta values},
\url{https://arxiv.org/abs/math/0206179}, which also records this
transcendence consequence for the divisor-sum series.
This is an exposition of the known resolution, not a new
algebraic-independence proof or a local formal verification.
The full question follows from the established external theorem.
\subsection{5) Completion estimate}
\textbf{COMPLETION: 100\%}
